\begin{apartat}{1}
Demostreu, a partir del principi de conservació de l'energia mecànica, que la velocitat d'escapament des d'un punt pròxim a la superfície d'un astre esfèric de massa $M$ i radi $R$ és $v_\mathrm{esc} = \sqrt{\dfrac{2GM}{R}}$.
\end{apartat}

\begin{apartat}{1}
Un objecte es llança verticalment des de la superfície de la Lluna amb una velocitat igual a la meitat de la velocitat d'escapament de la Lluna. Calculeu a quina altura màxima arribarà abans de tornar a caure.
\end{apartat}

\dades{%
  $G = 6,67\times 10^{-11}\ \mathrm{N\,m^2\,kg^{-2}}$\\
  $M_\mathrm{Lluna} = 7,35\times 10^{22}\ \mathrm{kg}$\\
  $R_\mathrm{Lluna} = 1\,737\ \mathrm{km}$}
